% titel: Blatt 1, A2: Summennotation

% Teil I: Werkzeuge, kein Lösungsweg.

Summen mit dem Summenzeichen schreiben, ausgewertete Werte für kleine
$n$ bestimmen und eine Summe mit ``\dots'' ausschreiben und berechnen.

\begin{werkzeugkasten}[Summenzeichen]
\begin{equation}
  \sum_{k=m}^{n} a_k = a_m + a_{m+1}+\dots+a_n .
\end{equation}
Den Summanden so umschreiben, dass ein gemeinsamer Laufindex entsteht. Bei
$2\cdot3+3\cdot4+\dots+n(n+1)$ ist der $k$-te Summand $k(k+1)$ mit
$k=2,\dots,n$.

\begin{itemize}
  \item Summen über Fakultäten: $(k+1)!=1\cdot2\cdots(k+1)$.
  \item Quadrat-Regel: $\sum_{k=0}^{n}k=\frac{n(n+1)}2$,
        $\sum_{i=1}^{n}(2i-1)=n^2$.
  \item Teleskopsumme: $\sum_{i=1}^{n}(a_i-a_{i-1})=a_n-a_0$.
\end{itemize}
\end{werkzeugkasten}

\begin{werkzeugkasten}[Teleskopprodukt/Teleskopsumme]
Summen der Form $\sum_{i=1}^{n}(i^2-(i-1)^2)$: Gesamt-Ausdruck
$i^2-(i-1)^2=2i-1$; oder direkt die Kette:
\begin{equation}
  \sum_{i=1}^{n}(i^2-(i-1)^2)=n^2-0^2=n^2 .
\end{equation}
das kürzt sich Schritt für Schritt weg.
\end{werkzeugkasten}

\begin{achtung}[Darauf achten]
\begin{itemize}
  \item Den Startwert korrekt ins Summenzeichen setzen (z.B. $k=2$ bei
        $2\cdot3+3\cdot4+\dots$).
  \item Fakultäten nur für kleine $n$ von Hand auswerten; große Summen per
        Formel oder mit ``\dots'' ausschreiben, nicht gliedweise rechnen.
  \item Bei der Teleskopsumme den letzten und ersten Term richtig
        identifizieren (mittlere heben sich weg).
\end{itemize}
\end{achtung}
